
\documentclass{article}
\usepackage{colortbl}
\usepackage{makecell}
\usepackage{multirow}
\usepackage{supertabular}

\begin{document}

\newcounter{utterance}

\centering \large Interaction Transcript for game `clean\_up', experiment `2\_hard\_5obj\_et', episode 2 with deepseek{-}v4{-}pro.
\vspace{24pt}

{ \footnotesize  \setcounter{utterance}{1}
\setlength{\tabcolsep}{0pt}
\begin{supertabular}{c@{$\;$}|p{.15\linewidth}@{}p{.15\linewidth}p{.15\linewidth}p{.15\linewidth}p{.15\linewidth}p{.15\linewidth}}
    \# & $\;$A & \multicolumn{4}{c}{Game Master} & $\;\:$B\\
    \hline

    \theutterance \stepcounter{utterance}  
    & & \multicolumn{4}{p{0.6\linewidth}}{
        \cellcolor[rgb]{0.9,0.9,0.9}{
            \makecell[{{p{\linewidth}}}]{
                \texttt{\tiny{[P1$\langle$GM]}}
                \texttt{Ma olen teie mängujuht ja te mängite ühismängu järgmise ruudustikuga mängulauana:} \\
\\ 
\texttt{```} \\
\texttt{ 1234567} \\
\texttt{╔═╤══╤══╗} \\
\texttt{║◌│◌◌│◌◌║ 1} \\
\texttt{║W└┬─┴──╢ 2} \\
\texttt{║◌◌│◌◌◌◌║ 3} \\
\texttt{╟─┬┘◌◌┌─╢ 4} \\
\texttt{║I│H◌◌│◌║ 5} \\
\texttt{╟─┼─┬─┤◌║ 6} \\
\texttt{║◌│C│T│◌║ 7} \\
\texttt{╚═╧═╧═╧═╝} \\
\\ 
\texttt{```} \\
\\ 
\texttt{* Ülemisel serval on x{-}koordinaadid, mis suurenevad paremale, ja paremal serval y{-}koordinaadid, mis suurenevad allapoole.} \\
\texttt{* Järgmised objektid on teie ruudustikule juhuslikult paigutatud: 'W', 'I', 'T', 'C', 'H'.} \\
\\ 
\\ 
\texttt{Teine mängija näeb mängulaua varianti, kus objektid on paigutatud erinevatesse juhuslikesse asukohtadesse. Te ei näe teise mängija mängulauda ja tema ei näe teie mängulauda.} \\
\\ 
\\ 
\texttt{**Eesmärk:**} \\
\\ 
\texttt{Mõlemad mängijad peavad oma taustal olevaid objekte liigutama nii, et identsed objektid satuksid samadesse koordinaatidesse. Te peate teise mängijaga suhtlema, et kokku leppida ühises eesmärgi konfiguratsioonis.} \\
\\ 
\\ 
\texttt{**Reeglid:**} \\
\\ 
\texttt{* Igal käigul saate saata täpselt ühe kahest järgmisest käsust:} \\
\texttt{ 1. `ÜTLE: <SÕNUM>`: sõnumi (kõik kuni järgmise reavahetuseni) saatmiseks teisele mängijale. Ma edastan selle teie partnerile.} \\
\texttt{2. `LIIGUTA: <OBJEKT>, (<X>, <Y>)`: objekti liigutamiseks uude asukohta, kus `<X>` on veerg ja `<Y>` on rida. Ma teatan teile, kas teie käik õnnestus või mitte.} \\
\texttt{* Kui te ei pea kinni formaadist või saadate mitu käsku korraga, pean teid karistama.} \\
\texttt{* Kui mõlemad mängijad koguvad rohkem kui 12 karistust, kaotate mõlemad mängu.} \\
\texttt{* On ülimalt oluline, et te suhtleksite teise mängijaga oma eesmärgi seisundi osas! *Ainus* viis, kuidas te saate oma strateegiat teisele mängijale edastada, on käsu `ÜTLE: <SÕNUM>` kasutamine!} \\
\\ 
\\ 
\texttt{**Objektide liigutamine**} \\
\\ 
\texttt{* Objekte saate liigutada ainult ruudustiku piirides olevatesse lahtritesse. Sihtlahter peab olema tühi, st see peab sisaldama ainult sümbolit '◌'.} \\
\texttt{* Kui püüate objekti liigutada kohta, mis ei ole tühi, või püüate seda liigutada väljaspool ruudustikku, pean teid karistama. Saate uue katse.} \\
\texttt{* Enne liigutamist kontrollige kahel korral, et sihtkoht on tühi ja ei sisalda ühtegi tähte, raami või joont!} \\
\\ 
\\ 
\texttt{**Mängu lõpp**} \\
\\ 
\texttt{Kui arvate, et olete saavutanud kõigi objektide joondamise eesmärgi, võite paluda teisel mängijal mäng lõpetada, saates `ÜTLE: valmis?`. Kui teine mängija palub teil mängu lõpetada ja te vastate `ÜTLE: valmis!`, siis mäng lõpeb.} \\
\\ 
\texttt{Mõlemad mängijad võidavad, kui mäng lõpeb 20 vooru jooksul, kus üks voor on määratletud kahe mängija poolt ühe kehtiva käsu saatmisena.} \\
\\ 
\texttt{**Punktiarvestus:**} \\
\\ 
\texttt{Mida lähemal on identsed objektid mõlemal mängulaual, seda rohkem punkte saate. Karistused vähendavad teie punkte. Kas suudate rekordit ületada?} \\
\\ 
\\ 
\texttt{Palun saada teisele mängijale sõnum, et mängu alustada!} \\
            }
        }
    }
    & \\ \\

    \theutterance \stepcounter{utterance}  
    & \multicolumn{4}{p{0.6\linewidth}}{
        \cellcolor[rgb]{0.9,0.9,0.9}{
            \makecell[{{p{\linewidth}}}]{
                \texttt{\tiny{[P1$\rangle$GM]}}
                \texttt{ÜTLE: Tere! Alustame mängu. Minu laual on objektid W (2,1), I (1,5), H (2,5), C (2,7) ja T (3,7). Palun kirjelda oma laudade seisundit, et saaksime leida ühised sihtkohad.} \\
            }
        }
    }
    & & \\ \\

    \theutterance \stepcounter{utterance}  
    & & & \multicolumn{2}{p{0.3\linewidth}}{
        \cellcolor[rgb]{0.9,0.9,0.9}{
            \makecell[{{p{\linewidth}}}]{
                \texttt{\tiny{[GM$|$GM]}}
                \texttt{ÜTLE: Tere! Alustame mängu. Minu laual on objektid W (2,1), I (1,5), H (2,5), C (2,7) ja T (3,7). Palun kirjelda oma laudade seisundit, et saaksime leida ühised sihtkohad.} \\
            }
        }
    }
    & & \\ \\

    \theutterance \stepcounter{utterance}  
    & & \multicolumn{4}{p{0.6\linewidth}}{
        \cellcolor[rgb]{0.9,0.9,0.9}{
            \makecell[{{p{\linewidth}}}]{
                \texttt{\tiny{[GM$\rangle$P2]}}
                \texttt{Ma olen teie mängujuht ja te mängite ühismängu järgmise ruudustikuga mängulauana:} \\
\\ 
\texttt{```} \\
\texttt{ 1234567} \\
\texttt{╔═╤══╤══╗} \\
\texttt{║◌│◌◌│◌H║ 1} \\
\texttt{║◌└┬─┴──╢ 2} \\
\texttt{║◌◌│◌◌◌◌║ 3} \\
\texttt{╟─┬┘◌I┌─╢ 4} \\
\texttt{║◌│WT◌│◌║ 5} \\
\texttt{╟─┼─┬─┤◌║ 6} \\
\texttt{║◌│C│◌│◌║ 7} \\
\texttt{╚═╧═╧═╧═╝} \\
\\ 
\texttt{```} \\
\\ 
\texttt{* Ülemisel serval on x{-}koordinaadid, mis suurenevad paremale, ja paremal serval y{-}koordinaadid, mis suurenevad allapoole.} \\
\texttt{* Järgmised objektid on teie ruudustikule juhuslikult paigutatud: 'W', 'I', 'T', 'C', 'H'.} \\
\\ 
\\ 
\texttt{Teine mängija näeb mängulaua varianti, kus objektid on paigutatud erinevatesse juhuslikesse asukohtadesse. Te ei näe teise mängija mängulauda ja tema ei näe teie mängulauda.} \\
\\ 
\\ 
\texttt{**Eesmärk:**} \\
\\ 
\texttt{Mõlemad mängijad peavad oma taustal olevaid objekte liigutama nii, et identsed objektid satuksid samadesse koordinaatidesse. Te peate teise mängijaga suhtlema, et kokku leppida ühises eesmärgi konfiguratsioonis.} \\
\\ 
\\ 
\texttt{**Reeglid:**} \\
\\ 
\texttt{* Igal käigul saate saata täpselt ühe kahest järgmisest käsust:} \\
\texttt{ 1. `ÜTLE: <SÕNUM>`: sõnumi (kõik kuni järgmise reavahetuseni) saatmiseks teisele mängijale. Ma edastan selle teie partnerile.} \\
\texttt{2. `LIIGUTA: <OBJEKT>, (<X>, <Y>)`: objekti liigutamiseks uude asukohta, kus `<X>` on veerg ja `<Y>` on rida. Ma teatan teile, kas teie käik õnnestus või mitte.} \\
\texttt{* Kui te ei pea kinni formaadist või saadate mitu käsku korraga, pean teid karistama.} \\
\texttt{* Kui mõlemad mängijad koguvad rohkem kui 12 karistust, kaotate mõlemad mängu.} \\
\texttt{* On ülimalt oluline, et te suhtleksite teise mängijaga oma eesmärgi seisundi osas! *Ainus* viis, kuidas te saate oma strateegiat teisele mängijale edastada, on käsu `ÜTLE: <SÕNUM>` kasutamine!} \\
\\ 
\\ 
\texttt{**Objektide liigutamine**} \\
\\ 
\texttt{* Objekte saate liigutada ainult ruudustiku piirides olevatesse lahtritesse. Sihtlahter peab olema tühi, st see peab sisaldama ainult sümbolit '◌'.} \\
\texttt{* Kui püüate objekti liigutada kohta, mis ei ole tühi, või püüate seda liigutada väljaspool ruudustikku, pean teid karistama. Saate uue katse.} \\
\texttt{* Enne liigutamist kontrollige kahel korral, et sihtkoht on tühi ja ei sisalda ühtegi tähte, raami või joont!} \\
\\ 
\\ 
\texttt{**Mängu lõpp**} \\
\\ 
\texttt{Kui arvate, et olete saavutanud kõigi objektide joondamise eesmärgi, võite paluda teisel mängijal mäng lõpetada, saates `ÜTLE: valmis?`. Kui teine mängija palub teil mängu lõpetada ja te vastate `ÜTLE: valmis!`, siis mäng lõpeb.} \\
\\ 
\texttt{Mõlemad mängijad võidavad, kui mäng lõpeb 20 vooru jooksul, kus üks voor on määratletud kahe mängija poolt ühe kehtiva käsu saatmisena.} \\
\\ 
\texttt{**Punktiarvestus:**} \\
\\ 
\texttt{Mida lähemal on identsed objektid mõlemal mängulaual, seda rohkem punkte saate. Karistused vähendavad teie punkte. Kas suudate rekordit ületada?} \\
\\ 
\texttt{Teine mängija alustas mängu selle sõnumi saatmisega:} \\
\\ 
\texttt{"Tere! Alustame mängu. Minu laual on objektid W (2,1), I (1,5), H (2,5), C (2,7) ja T (3,7). Palun kirjelda oma laudade seisundit, et saaksime leida ühised sihtkohad."} \\
\\ 
\texttt{Mis on sinu esimene käsk?} \\
            }
        }
    }
    & \\ \\

    \theutterance \stepcounter{utterance}  
    & & & \multicolumn{4}{p{0.6\linewidth}}{
        \cellcolor[rgb]{0.9,0.9,0.9}{
            \makecell[{{p{\linewidth}}}]{
                \texttt{\tiny{[GM$\langle$P2]}}
                \texttt{ÜTLE: Tere! Minu laual on objektid: W (4,5), I (4,4), T (3,5), C (3,7) ja H (7,1). Vaatame, kuidas need ühtida võiksid. Minu pakutud ühised sihtkohad: W (4,5), I (4,4), T (3,5), C (3,7), H (7,1) {-} need on juba minu laual olemas. Kas sul oleks võimalik oma objekte nendesse kohtadesse liigutada?} \\
            }
        }
    }
     \\ \\

    \theutterance \stepcounter{utterance}  
    & & & \multicolumn{2}{p{0.3\linewidth}}{
        \cellcolor[rgb]{0.9,0.9,0.9}{
            \makecell[{{p{\linewidth}}}]{
                \texttt{\tiny{[GM$|$GM]}}
                \texttt{ÜTLE: Tere! Minu laual on objektid: W (4,5), I (4,4), T (3,5), C (3,7) ja H (7,1). Vaatame, kuidas need ühtida võiksid. Minu pakutud ühised sihtkohad: W (4,5), I (4,4), T (3,5), C (3,7), H (7,1) {-} need on juba minu laual olemas. Kas sul oleks võimalik oma objekte nendesse kohtadesse liigutada?} \\
            }
        }
    }
    & & \\ \\

    \theutterance \stepcounter{utterance}  
    & & \multicolumn{4}{p{0.6\linewidth}}{
        \cellcolor[rgb]{0.9,0.9,0.9}{
            \makecell[{{p{\linewidth}}}]{
                \texttt{\tiny{[P1$\langle$GM]}}
                \texttt{Sinu sõnum on teisele mängijale edastatud.} \\
\\ 
\\ 
\texttt{Sa mängid praegu vooru 1 maksimaalsest 20 voorust.} \\
\texttt{Teil on kollektiivselt kogunenud 0 karistust 12 karistusest.} \\
\texttt{Teine mängija saatis selle sõnumi:} \\
\\ 
\\ 
\texttt{"Tere! Minu laual on objektid: W (4,5), I (4,4), T (3,5), C (3,7) ja H (7,1). Vaatame, kuidas need ühtida võiksid. Minu pakutud ühised sihtkohad: W (4,5), I (4,4), T (3,5), C (3,7), H (7,1) {-} need on juba minu laual olemas. Kas sul oleks võimalik oma objekte nendesse kohtadesse liigutada?"} \\
\\ 
\\ 
\texttt{Mis on sinu järgmine käsk?} \\
            }
        }
    }
    & \\ \\

    \theutterance \stepcounter{utterance}  
    & \multicolumn{4}{p{0.6\linewidth}}{
        \cellcolor[rgb]{0.9,0.9,0.9}{
            \makecell[{{p{\linewidth}}}]{
                \texttt{\tiny{[P1$\rangle$GM]}}
                \texttt{ÜTLE: Aitäh info eest! Kontrollime üle. Minu laual on W (2,1), I (1,5), H (2,5), C (2,7), T (3,7). Võrdleme sinu pakutud sihtkohtadega: W (4,5), I (4,4), T (3,5), C (3,7), H (7,1).} \\
\\ 
\texttt{Kontrollin oma laual, kas need kohad on tühjad:} \\
\texttt{{-} (4,5): ◌ (tühi)} \\
\texttt{{-} (4,4): ─ (joon, EI OLE tühi)} \\
\texttt{{-} (3,5): ◌ (tühi)} \\
\texttt{{-} (3,7): T (minu T on seal, aga sinu T on (3,5) {-} see on erinev koht)} \\
\texttt{{-} (7,1): ╗ (raam, EI OLE tühi)} \\
\\ 
\texttt{Probleemid: (4,4) ja (7,1) ei ole minu laual tühjad kohad. Kas saaksime leida alternatiivsed sihtkohad, mis on mõlema laual vabad? Palun pakku uued.} \\
            }
        }
    }
    & & \\ \\

    \theutterance \stepcounter{utterance}  
    & & & \multicolumn{2}{p{0.3\linewidth}}{
        \cellcolor[rgb]{0.9,0.9,0.9}{
            \makecell[{{p{\linewidth}}}]{
                \texttt{\tiny{[GM$|$GM]}}
                \texttt{ÜTLE: Aitäh info eest! Kontrollime üle. Minu laual on W (2,1), I (1,5), H (2,5), C (2,7), T (3,7). Võrdleme sinu pakutud sihtkohtadega: W (4,5), I (4,4), T (3,5), C (3,7), H (7,1).} \\
\\ 
\texttt{Kontrollin oma laual, kas need kohad on tühjad:} \\
\texttt{{-} (4,5): ◌ (tühi)} \\
\texttt{{-} (4,4): ─ (joon, EI OLE tühi)} \\
\texttt{{-} (3,5): ◌ (tühi)} \\
\texttt{{-} (3,7): T (minu T on seal, aga sinu T on (3,5) {-} see on erinev koht)} \\
\texttt{{-} (7,1): ╗ (raam, EI OLE tühi)} \\
\\ 
\texttt{Probleemid: (4,4) ja (7,1) ei ole minu laual tühjad kohad. Kas saaksime leida alternatiivsed sihtkohad, mis on mõlema laual vabad? Palun pakku uued.} \\
            }
        }
    }
    & & \\ \\

    \theutterance \stepcounter{utterance}  
    & & & \multicolumn{2}{p{0.3\linewidth}}{
        \cellcolor[rgb]{0.9,0.9,0.9}{
            \makecell[{{p{\linewidth}}}]{
                \texttt{\tiny{[GM$|$GM]}}
                \texttt{Invalid format: tail is not empty: '} \\
\\ 
\texttt{Kontrollin oma laual, kas need kohad on tühjad:} \\
\texttt{{-} (4,5): ◌ (tühi)} \\
\texttt{{-} (4,4): ─ (joon, EI OLE tühi)} \\
\texttt{{-} (3,5): ◌ (tühi)} \\
\texttt{{-} (3,7): T (minu T on seal, aga sinu T on (3,5) {-} see on erinev koht)} \\
\texttt{{-} (7,1): ╗ (raam, EI OLE tühi)} \\
\\ 
\texttt{Probleemid: (4,4) ja (7,1) ei ole minu laual tühjad kohad. Kas saaksime leida alternatiivsed sihtkohad, mis on mõlema laual vabad? Palun pakku uued.} \\
            }
        }
    }
    & & \\ \\

    \theutterance \stepcounter{utterance}  
    & & \multicolumn{4}{p{0.6\linewidth}}{
        \cellcolor[rgb]{0.9,0.9,0.9}{
            \makecell[{{p{\linewidth}}}]{
                \texttt{\tiny{[P1$\langle$GM]}}
                \texttt{Karistus: Teie sõnum ei tohi sisaldada midagi pärast käsku!} \\
\\ 
\texttt{Veendu, et sinu vastus sisaldab ainult kas `ÜTLE: <SÕNUM>` või `LIIGUTA: <OBJEKT>, (<X>, <Y>)` ja mitte midagi muud!} \\
\\ 
\texttt{Teil on kollektiivselt kogunenud 1 karistust 12 karistusest.} \\
\texttt{Palun proovi uuesti!} \\
            }
        }
    }
    & \\ \\

    \theutterance \stepcounter{utterance}  
    & \multicolumn{4}{p{0.6\linewidth}}{
        \cellcolor[rgb]{0.9,0.9,0.9}{
            \makecell[{{p{\linewidth}}}]{
                \texttt{\tiny{[P1$\rangle$GM]}}
                \texttt{ÜTLE: Aitäh info eest! Kontrollisin sinu pakutud sihtkohti. Minu laual ei ole (4,4) ja (7,1) tühjad, need sisaldavad jooni. Kas saaksime leida alternatiivsed sihtkohad, mis on mõlema laual vabad? Palun paku uued võimalikud asukohad.} \\
            }
        }
    }
    & & \\ \\

    \theutterance \stepcounter{utterance}  
    & & & \multicolumn{2}{p{0.3\linewidth}}{
        \cellcolor[rgb]{0.9,0.9,0.9}{
            \makecell[{{p{\linewidth}}}]{
                \texttt{\tiny{[GM$|$GM]}}
                \texttt{ÜTLE: Aitäh info eest! Kontrollisin sinu pakutud sihtkohti. Minu laual ei ole (4,4) ja (7,1) tühjad, need sisaldavad jooni. Kas saaksime leida alternatiivsed sihtkohad, mis on mõlema laual vabad? Palun paku uued võimalikud asukohad.} \\
            }
        }
    }
    & & \\ \\

    \theutterance \stepcounter{utterance}  
    & & \multicolumn{4}{p{0.6\linewidth}}{
        \cellcolor[rgb]{0.9,0.9,0.9}{
            \makecell[{{p{\linewidth}}}]{
                \texttt{\tiny{[GM$\rangle$P2]}}
                \texttt{Sinu sõnum on teisele mängijale edastatud.} \\
\\ 
\\ 
\texttt{Sa mängid praegu vooru 2 maksimaalsest 20 voorust.} \\
\texttt{Teil on kollektiivselt kogunenud 1 karistust 12 karistusest.} \\
\texttt{Teine mängija saatis selle sõnumi:} \\
\\ 
\\ 
\texttt{"Aitäh info eest! Kontrollisin sinu pakutud sihtkohti. Minu laual ei ole (4,4) ja (7,1) tühjad, need sisaldavad jooni. Kas saaksime leida alternatiivsed sihtkohad, mis on mõlema laual vabad? Palun paku uued võimalikud asukohad."} \\
\\ 
\\ 
\texttt{Mis on sinu järgmine käsk?} \\
            }
        }
    }
    & \\ \\

    \theutterance \stepcounter{utterance}  
    & & & \multicolumn{4}{p{0.6\linewidth}}{
        \cellcolor[rgb]{0.9,0.9,0.9}{
            \makecell[{{p{\linewidth}}}]{
                \texttt{\tiny{[GM$\langle$P2]}}
                \texttt{ÜTLE: Vabandust, mul polnud seda infot. Vaatame minu laual vabu kohti: (1,1), (2,1), (3,1), (4,1), (5,1), (6,1), (1,2), (2,2), (3,2), (6,2), (1,3), (2,3), (5,3), (6,3), (7,3), (1,4), (3,4), (6,4), (7,4), (1,5), (2,5), (5,5), (6,5), (7,5), (1,6), (2,6), (4,6), (5,6), (7,6), (1,7), (2,7), (4,7), (5,7), (6,7), (7,7). Kas sinu laual on kõik need tühjad? Kui jah, siis pakun uued sihtkohad: W (2,5) {-} see sul juba on, I (6,5), T (5,5), C (4,7), H (1,5) {-} H sul juba on. Kinnita, kas need sobivad.} \\
            }
        }
    }
     \\ \\

    \theutterance \stepcounter{utterance}  
    & & & \multicolumn{2}{p{0.3\linewidth}}{
        \cellcolor[rgb]{0.9,0.9,0.9}{
            \makecell[{{p{\linewidth}}}]{
                \texttt{\tiny{[GM$|$GM]}}
                \texttt{ÜTLE: Vabandust, mul polnud seda infot. Vaatame minu laual vabu kohti: (1,1), (2,1), (3,1), (4,1), (5,1), (6,1), (1,2), (2,2), (3,2), (6,2), (1,3), (2,3), (5,3), (6,3), (7,3), (1,4), (3,4), (6,4), (7,4), (1,5), (2,5), (5,5), (6,5), (7,5), (1,6), (2,6), (4,6), (5,6), (7,6), (1,7), (2,7), (4,7), (5,7), (6,7), (7,7). Kas sinu laual on kõik need tühjad? Kui jah, siis pakun uued sihtkohad: W (2,5) {-} see sul juba on, I (6,5), T (5,5), C (4,7), H (1,5) {-} H sul juba on. Kinnita, kas need sobivad.} \\
            }
        }
    }
    & & \\ \\

    \theutterance \stepcounter{utterance}  
    & & \multicolumn{4}{p{0.6\linewidth}}{
        \cellcolor[rgb]{0.9,0.9,0.9}{
            \makecell[{{p{\linewidth}}}]{
                \texttt{\tiny{[P1$\langle$GM]}}
                \texttt{Sinu sõnum on teisele mängijale edastatud.} \\
\\ 
\\ 
\texttt{Sa mängid praegu vooru 2 maksimaalsest 20 voorust.} \\
\texttt{Teil on kollektiivselt kogunenud 1 karistust 12 karistusest.} \\
\texttt{Teine mängija saatis selle sõnumi:} \\
\\ 
\\ 
\texttt{"Vabandust, mul polnud seda infot. Vaatame minu laual vabu kohti: (1,1), (2,1), (3,1), (4,1), (5,1), (6,1), (1,2), (2,2), (3,2), (6,2), (1,3), (2,3), (5,3), (6,3), (7,3), (1,4), (3,4), (6,4), (7,4), (1,5), (2,5), (5,5), (6,5), (7,5), (1,6), (2,6), (4,6), (5,6), (7,6), (1,7), (2,7), (4,7), (5,7), (6,7), (7,7). Kas sinu laual on kõik need tühjad? Kui jah, siis pakun uued sihtkohad: W (2,5) {-} see sul juba on, I (6,5), T (5,5), C (4,7), H (1,5) {-} H sul juba on. Kinnita, kas need sobivad."} \\
\\ 
\\ 
\texttt{Mis on sinu järgmine käsk?} \\
            }
        }
    }
    & \\ \\

    \theutterance \stepcounter{utterance}  
    & \multicolumn{4}{p{0.6\linewidth}}{
        \cellcolor[rgb]{0.9,0.9,0.9}{
            \makecell[{{p{\linewidth}}}]{
                \texttt{\tiny{[P1$\rangle$GM]}}
                \texttt{ÜTLE: Kontrollisin sinu uued pakutud sihtkohad. Minu laual: (2,5) on H, mitte W {-} seal on objekt. (1,5) on I, mitte H. (6,5) on ◌ (tühi), (5,5) on ◌ (tühi), (4,7) on ◌ (tühi). Minu laual on W (2,1), I (1,5), H (2,5), C (2,7), T (3,7). Kas saaksime kasutada hoopis neid sihtkohti: W (4,5), I (6,5), H (5,5), C (4,7), T (3,5)? Minu laual on (4,5), (6,5), (5,5), (4,7), (3,5) kõik tühjad. Kinnita, kas need sobivad sinu lauale.} \\
            }
        }
    }
    & & \\ \\

    \theutterance \stepcounter{utterance}  
    & & & \multicolumn{2}{p{0.3\linewidth}}{
        \cellcolor[rgb]{0.9,0.9,0.9}{
            \makecell[{{p{\linewidth}}}]{
                \texttt{\tiny{[GM$|$GM]}}
                \texttt{ÜTLE: Kontrollisin sinu uued pakutud sihtkohad. Minu laual: (2,5) on H, mitte W {-} seal on objekt. (1,5) on I, mitte H. (6,5) on ◌ (tühi), (5,5) on ◌ (tühi), (4,7) on ◌ (tühi). Minu laual on W (2,1), I (1,5), H (2,5), C (2,7), T (3,7). Kas saaksime kasutada hoopis neid sihtkohti: W (4,5), I (6,5), H (5,5), C (4,7), T (3,5)? Minu laual on (4,5), (6,5), (5,5), (4,7), (3,5) kõik tühjad. Kinnita, kas need sobivad sinu lauale.} \\
            }
        }
    }
    & & \\ \\

    \theutterance \stepcounter{utterance}  
    & & \multicolumn{4}{p{0.6\linewidth}}{
        \cellcolor[rgb]{0.9,0.9,0.9}{
            \makecell[{{p{\linewidth}}}]{
                \texttt{\tiny{[GM$\rangle$P2]}}
                \texttt{Sinu sõnum on teisele mängijale edastatud.} \\
\\ 
\\ 
\texttt{Sa mängid praegu vooru 3 maksimaalsest 20 voorust.} \\
\texttt{Teil on kollektiivselt kogunenud 1 karistust 12 karistusest.} \\
\texttt{Teine mängija saatis selle sõnumi:} \\
\\ 
\\ 
\texttt{"Kontrollisin sinu uued pakutud sihtkohad. Minu laual: (2,5) on H, mitte W {-} seal on objekt. (1,5) on I, mitte H. (6,5) on ◌ (tühi), (5,5) on ◌ (tühi), (4,7) on ◌ (tühi). Minu laual on W (2,1), I (1,5), H (2,5), C (2,7), T (3,7). Kas saaksime kasutada hoopis neid sihtkohti: W (4,5), I (6,5), H (5,5), C (4,7), T (3,5)? Minu laual on (4,5), (6,5), (5,5), (4,7), (3,5) kõik tühjad. Kinnita, kas need sobivad sinu lauale."} \\
\\ 
\\ 
\texttt{Mis on sinu järgmine käsk?} \\
            }
        }
    }
    & \\ \\

    \theutterance \stepcounter{utterance}  
    & & & \multicolumn{4}{p{0.6\linewidth}}{
        \cellcolor[rgb]{0.9,0.9,0.9}{
            \makecell[{{p{\linewidth}}}]{
                \texttt{\tiny{[GM$\langle$P2]}}
                \texttt{ÜTLE: Kontrollisin sinu uut pakkumist minu laua peal. Minu laual: (4,5) on W {-} juba õiges kohas, (6,5) on tühi {-} sobib I{-}le, (5,5) on tühi {-} sobib H{-}le, (4,7) on tühi {-} sobib C{-}le, (3,5) on T {-} juba õiges kohas. Suurepärane, need sobivad ideaalselt! Minu objektid W ja T on juba õigetes kohtades. Hakkan nüüd liigutama I, H ja C{-}d nendesse sihtkohtadesse. Sina peaksid liigutama W (2,1) {-}> (4,5), I (1,5) {-}> (6,5), H (2,5) {-}> (5,5), C (2,7) {-}> (4,7) ja T (3,7) {-}> (3,5). Alustame liigutamist!} \\
            }
        }
    }
     \\ \\

    \theutterance \stepcounter{utterance}  
    & & & \multicolumn{2}{p{0.3\linewidth}}{
        \cellcolor[rgb]{0.9,0.9,0.9}{
            \makecell[{{p{\linewidth}}}]{
                \texttt{\tiny{[GM$|$GM]}}
                \texttt{ÜTLE: Kontrollisin sinu uut pakkumist minu laua peal. Minu laual: (4,5) on W {-} juba õiges kohas, (6,5) on tühi {-} sobib I{-}le, (5,5) on tühi {-} sobib H{-}le, (4,7) on tühi {-} sobib C{-}le, (3,5) on T {-} juba õiges kohas. Suurepärane, need sobivad ideaalselt! Minu objektid W ja T on juba õigetes kohtades. Hakkan nüüd liigutama I, H ja C{-}d nendesse sihtkohtadesse. Sina peaksid liigutama W (2,1) {-}> (4,5), I (1,5) {-}> (6,5), H (2,5) {-}> (5,5), C (2,7) {-}> (4,7) ja T (3,7) {-}> (3,5). Alustame liigutamist!} \\
            }
        }
    }
    & & \\ \\

    \theutterance \stepcounter{utterance}  
    & & \multicolumn{4}{p{0.6\linewidth}}{
        \cellcolor[rgb]{0.9,0.9,0.9}{
            \makecell[{{p{\linewidth}}}]{
                \texttt{\tiny{[P1$\langle$GM]}}
                \texttt{Sinu sõnum on teisele mängijale edastatud.} \\
\\ 
\\ 
\texttt{Sa mängid praegu vooru 3 maksimaalsest 20 voorust.} \\
\texttt{Teil on kollektiivselt kogunenud 1 karistust 12 karistusest.} \\
\texttt{Teine mängija saatis selle sõnumi:} \\
\\ 
\\ 
\texttt{"Kontrollisin sinu uut pakkumist minu laua peal. Minu laual: (4,5) on W {-} juba õiges kohas, (6,5) on tühi {-} sobib I{-}le, (5,5) on tühi {-} sobib H{-}le, (4,7) on tühi {-} sobib C{-}le, (3,5) on T {-} juba õiges kohas. Suurepärane, need sobivad ideaalselt! Minu objektid W ja T on juba õigetes kohtades. Hakkan nüüd liigutama I, H ja C{-}d nendesse sihtkohtadesse. Sina peaksid liigutama W (2,1) {-}> (4,5), I (1,5) {-}> (6,5), H (2,5) {-}> (5,5), C (2,7) {-}> (4,7) ja T (3,7) {-}> (3,5). Alustame liigutamist!"} \\
\\ 
\\ 
\texttt{Mis on sinu järgmine käsk?} \\
            }
        }
    }
    & \\ \\

    \theutterance \stepcounter{utterance}  
    & \multicolumn{4}{p{0.6\linewidth}}{
        \cellcolor[rgb]{0.9,0.9,0.9}{
            \makecell[{{p{\linewidth}}}]{
                \texttt{\tiny{[P1$\rangle$GM]}}
                \texttt{LIIGUTA: W, (4, 5)} \\
            }
        }
    }
    & & \\ \\

    \theutterance \stepcounter{utterance}  
    & & & \multicolumn{2}{p{0.3\linewidth}}{
        \cellcolor[rgb]{0.9,0.9,0.9}{
            \makecell[{{p{\linewidth}}}]{
                \texttt{\tiny{[GM$|$GM]}}
                \texttt{LIIGUTA: W, (4, 5)} \\
            }
        }
    }
    & & \\ \\

    \theutterance \stepcounter{utterance}  
    & & & \multicolumn{2}{p{0.3\linewidth}}{
        \cellcolor[rgb]{0.9,0.9,0.9}{
            \makecell[{{p{\linewidth}}}]{
                \texttt{\tiny{[GM$|$GM]}}
                \texttt{W moved to (4, 5)} \\
            }
        }
    }
    & & \\ \\

    \theutterance \stepcounter{utterance}  
    & & \multicolumn{4}{p{0.6\linewidth}}{
        \cellcolor[rgb]{0.9,0.9,0.9}{
            \makecell[{{p{\linewidth}}}]{
                \texttt{\tiny{[GM$\rangle$P2]}}
                \texttt{Sinu sõnum on teisele mängijale edastatud.} \\
\\ 
\\ 
\texttt{Sa mängid praegu vooru 4 maksimaalsest 20 voorust.} \\
\texttt{Teil on kollektiivselt kogunenud 1 karistust 12 karistusest.} \\
\texttt{Teine mängija liigutas objekti oma ruudustikul.} \\
\\ 
\\ 
\texttt{Mis on sinu järgmine käsk?} \\
            }
        }
    }
    & \\ \\

    \theutterance \stepcounter{utterance}  
    & & & \multicolumn{4}{p{0.6\linewidth}}{
        \cellcolor[rgb]{0.9,0.9,0.9}{
            \makecell[{{p{\linewidth}}}]{
                \texttt{\tiny{[GM$\langle$P2]}}
                \texttt{LIIGUTA: I, (4, 4)} \\
            }
        }
    }
     \\ \\

    \theutterance \stepcounter{utterance}  
    & & & \multicolumn{2}{p{0.3\linewidth}}{
        \cellcolor[rgb]{0.9,0.9,0.9}{
            \makecell[{{p{\linewidth}}}]{
                \texttt{\tiny{[GM$|$GM]}}
                \texttt{LIIGUTA: I, (4, 4)} \\
            }
        }
    }
    & & \\ \\

    \theutterance \stepcounter{utterance}  
    & & & \multicolumn{2}{p{0.3\linewidth}}{
        \cellcolor[rgb]{0.9,0.9,0.9}{
            \makecell[{{p{\linewidth}}}]{
                \texttt{\tiny{[GM$|$GM]}}
                \texttt{I moved to (4, 4)} \\
            }
        }
    }
    & & \\ \\

    \theutterance \stepcounter{utterance}  
    & & \multicolumn{4}{p{0.6\linewidth}}{
        \cellcolor[rgb]{0.9,0.9,0.9}{
            \makecell[{{p{\linewidth}}}]{
                \texttt{\tiny{[P1$\langle$GM]}}
                \texttt{Edukalt teisaldatud 'W' asukohta (4,5). Teie uuendatud ruudustik näeb välja selline:} \\
\texttt{```} \\
\texttt{ 1234567} \\
\texttt{╔═╤══╤══╗} \\
\texttt{║◌│◌◌│◌◌║ 1} \\
\texttt{║◌└┬─┴──╢ 2} \\
\texttt{║◌◌│◌◌◌◌║ 3} \\
\texttt{╟─┬┘◌◌┌─╢ 4} \\
\texttt{║I│HW◌│◌║ 5} \\
\texttt{╟─┼─┬─┤◌║ 6} \\
\texttt{║◌│C│T│◌║ 7} \\
\texttt{╚═╧═╧═╧═╝} \\
\texttt{```} \\
\texttt{Sa mängid praegu vooru 4 maksimaalsest 20 voorust.} \\
\texttt{Teil on kollektiivselt kogunenud 1 karistust 12 karistusest.} \\
\texttt{Teine mängija liigutas objekti oma ruudustikul.} \\
\\ 
\\ 
\texttt{Mis on sinu järgmine käsk?} \\
            }
        }
    }
    & \\ \\

    \theutterance \stepcounter{utterance}  
    & \multicolumn{4}{p{0.6\linewidth}}{
        \cellcolor[rgb]{0.9,0.9,0.9}{
            \makecell[{{p{\linewidth}}}]{
                \texttt{\tiny{[P1$\rangle$GM]}}
                \texttt{LIIGUTA: I, (6, 5)} \\
            }
        }
    }
    & & \\ \\

    \theutterance \stepcounter{utterance}  
    & & & \multicolumn{2}{p{0.3\linewidth}}{
        \cellcolor[rgb]{0.9,0.9,0.9}{
            \makecell[{{p{\linewidth}}}]{
                \texttt{\tiny{[GM$|$GM]}}
                \texttt{LIIGUTA: I, (6, 5)} \\
            }
        }
    }
    & & \\ \\

    \theutterance \stepcounter{utterance}  
    & & & \multicolumn{2}{p{0.3\linewidth}}{
        \cellcolor[rgb]{0.9,0.9,0.9}{
            \makecell[{{p{\linewidth}}}]{
                \texttt{\tiny{[GM$|$GM]}}
                \texttt{Karistus: (6,5) ei ole tühi, vaid sisaldab objekti '│'.} \\
\\ 
\texttt{Teil on kollektiivselt kogunenud 2 karistust 12 karistusest.} \\
\texttt{Palun proovi uuesti!} \\
            }
        }
    }
    & & \\ \\

    \theutterance \stepcounter{utterance}  
    & & \multicolumn{4}{p{0.6\linewidth}}{
        \cellcolor[rgb]{0.9,0.9,0.9}{
            \makecell[{{p{\linewidth}}}]{
                \texttt{\tiny{[P1$\langle$GM]}}
                \texttt{Karistus: (6,5) ei ole tühi, vaid sisaldab objekti '│'.} \\
\\ 
\texttt{Teil on kollektiivselt kogunenud 2 karistust 12 karistusest.} \\
\texttt{Palun proovi uuesti!} \\
            }
        }
    }
    & \\ \\

    \theutterance \stepcounter{utterance}  
    & \multicolumn{4}{p{0.6\linewidth}}{
        \cellcolor[rgb]{0.9,0.9,0.9}{
            \makecell[{{p{\linewidth}}}]{
                \texttt{\tiny{[P1$\rangle$GM]}}
                \texttt{ÜTLE: Minu laual (6,5) ei ole tühi, seal on '│'. Kas sa näed oma laual, et (6,5) on tühi? Võib{-}olla peame I jaoks uue sihtkoha leidma. Mis on sinu laual tühjad kohad veergudes 5{-}7 ja ridades 5{-}7?} \\
            }
        }
    }
    & & \\ \\

    \theutterance \stepcounter{utterance}  
    & & & \multicolumn{2}{p{0.3\linewidth}}{
        \cellcolor[rgb]{0.9,0.9,0.9}{
            \makecell[{{p{\linewidth}}}]{
                \texttt{\tiny{[GM$|$GM]}}
                \texttt{ÜTLE: Minu laual (6,5) ei ole tühi, seal on '│'. Kas sa näed oma laual, et (6,5) on tühi? Võib{-}olla peame I jaoks uue sihtkoha leidma. Mis on sinu laual tühjad kohad veergudes 5{-}7 ja ridades 5{-}7?} \\
            }
        }
    }
    & & \\ \\

    \theutterance \stepcounter{utterance}  
    & & \multicolumn{4}{p{0.6\linewidth}}{
        \cellcolor[rgb]{0.9,0.9,0.9}{
            \makecell[{{p{\linewidth}}}]{
                \texttt{\tiny{[GM$\rangle$P2]}}
                \texttt{Edukalt teisaldatud 'I' asukohta (4,4). Teie uuendatud ruudustik näeb välja selline:} \\
\texttt{```} \\
\texttt{ 1234567} \\
\texttt{╔═╤══╤══╗} \\
\texttt{║◌│◌◌│◌H║ 1} \\
\texttt{║◌└┬─┴──╢ 2} \\
\texttt{║◌◌│◌◌◌◌║ 3} \\
\texttt{╟─┬┘I◌┌─╢ 4} \\
\texttt{║◌│WT◌│◌║ 5} \\
\texttt{╟─┼─┬─┤◌║ 6} \\
\texttt{║◌│C│◌│◌║ 7} \\
\texttt{╚═╧═╧═╧═╝} \\
\texttt{```} \\
\texttt{Sa mängid praegu vooru 5 maksimaalsest 20 voorust.} \\
\texttt{Teil on kollektiivselt kogunenud 2 karistust 12 karistusest.} \\
\texttt{Teine mängija saatis selle sõnumi:} \\
\\ 
\\ 
\texttt{"Minu laual (6,5) ei ole tühi, seal on '│'. Kas sa näed oma laual, et (6,5) on tühi? Võib{-}olla peame I jaoks uue sihtkoha leidma. Mis on sinu laual tühjad kohad veergudes 5{-}7 ja ridades 5{-}7?"} \\
\\ 
\\ 
\texttt{Mis on sinu järgmine käsk?} \\
            }
        }
    }
    & \\ \\

    \theutterance \stepcounter{utterance}  
    & & & \multicolumn{4}{p{0.6\linewidth}}{
        \cellcolor[rgb]{0.9,0.9,0.9}{
            \makecell[{{p{\linewidth}}}]{
                \texttt{\tiny{[GM$\langle$P2]}}
                \texttt{Vabandust, tegin vea I liigutamisega. Ma pidin I liigutama (6,5), mitte (4,4). Aga nüüd näen, et minu laual (6,5) on tõepoolest tühi (◌). Aga teine mängija ütleb, et tema laual on seal '│'. Seega peame leidma I{-}le uue sihtkoha.} \\
\\ 
\texttt{Minu laual on veergudes 5{-}7 ja ridades 5{-}7 järgmised tühjad kohad:} \\
\texttt{(5,5), (6,5), (7,5) {-} kõik tühjad} \\
\texttt{(5,6), (7,6) {-} tühjad, aga (6,6) on '┤' ja (4,6) on tühi} \\
\texttt{(4,7), (5,7), (6,7), (7,7) {-} tühjad, aga (2,7) on C juba õiges kohas? Ei, C on (3,7).} \\
\\ 
\texttt{Sihtkohad pidid olema: W (4,5), I (4,4) {-} see oli minu viga, pidi olema (6,5), H (5,5), C (4,7), T (3,5).} \\
\\ 
\texttt{Aga kui (6,5) ei sobi I{-}le, siis peame leidma uue koha. Kas (7,5) sobiks? See on tühi nii minu kui loodetavasti ka sinu laual.} \\
\\ 
\texttt{ÜTLE: Mul on kahju segaduse pärast. Liigutasin I kogemata (4,4) asemel (4,4) {-} tegelikult pidin liigutama (6,5), aga see ei sobi sulle. Kas (7,5) on sinu laual tühi? See võiks olla I uus sihtkoht. Minu laual on (7,5) tühi. Kinnitage, kas see sobib, siis liigutan I sinna.} \\
            }
        }
    }
     \\ \\

    \theutterance \stepcounter{utterance}  
    & & & \multicolumn{2}{p{0.3\linewidth}}{
        \cellcolor[rgb]{0.9,0.9,0.9}{
            \makecell[{{p{\linewidth}}}]{
                \texttt{\tiny{[GM$|$GM]}}
                \texttt{Vabandust, tegin vea I liigutamisega. Ma pidin I liigutama (6,5), mitte (4,4). Aga nüüd näen, et minu laual (6,5) on tõepoolest tühi (◌). Aga teine mängija ütleb, et tema laual on seal '│'. Seega peame leidma I{-}le uue sihtkoha.} \\
\\ 
\texttt{Minu laual on veergudes 5{-}7 ja ridades 5{-}7 järgmised tühjad kohad:} \\
\texttt{(5,5), (6,5), (7,5) {-} kõik tühjad} \\
\texttt{(5,6), (7,6) {-} tühjad, aga (6,6) on '┤' ja (4,6) on tühi} \\
\texttt{(4,7), (5,7), (6,7), (7,7) {-} tühjad, aga (2,7) on C juba õiges kohas? Ei, C on (3,7).} \\
\\ 
\texttt{Sihtkohad pidid olema: W (4,5), I (4,4) {-} see oli minu viga, pidi olema (6,5), H (5,5), C (4,7), T (3,5).} \\
\\ 
\texttt{Aga kui (6,5) ei sobi I{-}le, siis peame leidma uue koha. Kas (7,5) sobiks? See on tühi nii minu kui loodetavasti ka sinu laual.} \\
\\ 
\texttt{ÜTLE: Mul on kahju segaduse pärast. Liigutasin I kogemata (4,4) asemel (4,4) {-} tegelikult pidin liigutama (6,5), aga see ei sobi sulle. Kas (7,5) on sinu laual tühi? See võiks olla I uus sihtkoht. Minu laual on (7,5) tühi. Kinnitage, kas see sobib, siis liigutan I sinna.} \\
            }
        }
    }
    & & \\ \\

    \theutterance \stepcounter{utterance}  
    & & & \multicolumn{2}{p{0.3\linewidth}}{
        \cellcolor[rgb]{0.9,0.9,0.9}{
            \makecell[{{p{\linewidth}}}]{
                \texttt{\tiny{[GM$|$GM]}}
                \texttt{Invalid format: head is not empty: 'Vabandust, tegin vea I liigutamisega. Ma pidin I liigutama (6,5), mitte (4,4). Aga nüüd näen, et minu laual (6,5) on tõepoolest tühi (◌). Aga teine mängija ütleb, et tema laual on seal '│'. Seega peame leidma I{-}le uue sihtkoha.} \\
\\ 
\texttt{Minu laual on veergudes 5{-}7 ja ridades 5{-}7 järgmised tühjad kohad:} \\
\texttt{(5,5), (6,5), (7,5) {-} kõik tühjad} \\
\texttt{(5,6), (7,6) {-} tühjad, aga (6,6) on '┤' ja (4,6) on tühi} \\
\texttt{(4,7), (5,7), (6,7), (7,7) {-} tühjad, aga (2,7) on C juba õiges kohas? Ei, C on (3,7).} \\
\\ 
\texttt{Sihtkohad pidid olema: W (4,5), I (4,4) {-} see oli minu viga, pidi olema (6,5), H (5,5), C (4,7), T (3,5).} \\
\\ 
\texttt{Aga kui (6,5) ei sobi I{-}le, siis peame leidma uue koha. Kas (7,5) sobiks? See on tühi nii minu kui loodetavasti ka sinu laual.} \\
\\ 
\texttt{'} \\
            }
        }
    }
    & & \\ \\

    \theutterance \stepcounter{utterance}  
    & & \multicolumn{4}{p{0.6\linewidth}}{
        \cellcolor[rgb]{0.9,0.9,0.9}{
            \makecell[{{p{\linewidth}}}]{
                \texttt{\tiny{[GM$\rangle$P2]}}
                \texttt{Karistus: Teie sõnum ei tohi sisaldada midagi enne käsku!} \\
\\ 
\texttt{Veendu, et sinu vastus sisaldab ainult kas `ÜTLE: <SÕNUM>` või `LIIGUTA: <OBJEKT>, (<X>, <Y>)` ja mitte midagi muud!} \\
\\ 
\texttt{Teil on kollektiivselt kogunenud 3 karistust 12 karistusest.} \\
\texttt{Palun proovi uuesti!} \\
            }
        }
    }
    & \\ \\

    \theutterance \stepcounter{utterance}  
    & & & \multicolumn{4}{p{0.6\linewidth}}{
        \cellcolor[rgb]{0.9,0.9,0.9}{
            \makecell[{{p{\linewidth}}}]{
                \texttt{\tiny{[GM$\langle$P2]}}
                \texttt{ÜTLE: Mul on kahju segaduse pärast. Liigutasin I kogemata (4,4) asemel (4,4) {-} tegelikult pidin liigutama (6,5), aga see ei sobi sulle. Kas (7,5) on sinu laual tühi? See võiks olla I uus sihtkoht. Minu laual on (7,5) tühi. Kinnitage, kas see sobib, siis liigutan I sinna.} \\
            }
        }
    }
     \\ \\

    \theutterance \stepcounter{utterance}  
    & & & \multicolumn{2}{p{0.3\linewidth}}{
        \cellcolor[rgb]{0.9,0.9,0.9}{
            \makecell[{{p{\linewidth}}}]{
                \texttt{\tiny{[GM$|$GM]}}
                \texttt{ÜTLE: Mul on kahju segaduse pärast. Liigutasin I kogemata (4,4) asemel (4,4) {-} tegelikult pidin liigutama (6,5), aga see ei sobi sulle. Kas (7,5) on sinu laual tühi? See võiks olla I uus sihtkoht. Minu laual on (7,5) tühi. Kinnitage, kas see sobib, siis liigutan I sinna.} \\
            }
        }
    }
    & & \\ \\

    \theutterance \stepcounter{utterance}  
    & & \multicolumn{4}{p{0.6\linewidth}}{
        \cellcolor[rgb]{0.9,0.9,0.9}{
            \makecell[{{p{\linewidth}}}]{
                \texttt{\tiny{[P1$\langle$GM]}}
                \texttt{Sinu sõnum on teisele mängijale edastatud.} \\
\\ 
\\ 
\texttt{Sa mängid praegu vooru 5 maksimaalsest 20 voorust.} \\
\texttt{Teil on kollektiivselt kogunenud 3 karistust 12 karistusest.} \\
\texttt{Teine mängija saatis selle sõnumi:} \\
\\ 
\\ 
\texttt{"Mul on kahju segaduse pärast. Liigutasin I kogemata (4,4) asemel (4,4) {-} tegelikult pidin liigutama (6,5), aga see ei sobi sulle. Kas (7,5) on sinu laual tühi? See võiks olla I uus sihtkoht. Minu laual on (7,5) tühi. Kinnitage, kas see sobib, siis liigutan I sinna."} \\
\\ 
\\ 
\texttt{Mis on sinu järgmine käsk?} \\
            }
        }
    }
    & \\ \\

    \theutterance \stepcounter{utterance}  
    & \multicolumn{4}{p{0.6\linewidth}}{
        \cellcolor[rgb]{0.9,0.9,0.9}{
            \makecell[{{p{\linewidth}}}]{
                \texttt{\tiny{[P1$\rangle$GM]}}
                \texttt{ÜTLE: Kontrollisin oma laual (7,5) {-} see on tühi ◌. Sobib I{-}le. Liigutan I sinna ja sina liiguta ka oma I (7,5){-}sse. Alustan liigutamist.} \\
            }
        }
    }
    & & \\ \\

    \theutterance \stepcounter{utterance}  
    & & & \multicolumn{2}{p{0.3\linewidth}}{
        \cellcolor[rgb]{0.9,0.9,0.9}{
            \makecell[{{p{\linewidth}}}]{
                \texttt{\tiny{[GM$|$GM]}}
                \texttt{ÜTLE: Kontrollisin oma laual (7,5) {-} see on tühi ◌. Sobib I{-}le. Liigutan I sinna ja sina liiguta ka oma I (7,5){-}sse. Alustan liigutamist.} \\
            }
        }
    }
    & & \\ \\

    \theutterance \stepcounter{utterance}  
    & & \multicolumn{4}{p{0.6\linewidth}}{
        \cellcolor[rgb]{0.9,0.9,0.9}{
            \makecell[{{p{\linewidth}}}]{
                \texttt{\tiny{[GM$\rangle$P2]}}
                \texttt{Sinu sõnum on teisele mängijale edastatud.} \\
\\ 
\\ 
\texttt{Sa mängid praegu vooru 6 maksimaalsest 20 voorust.} \\
\texttt{Teil on kollektiivselt kogunenud 3 karistust 12 karistusest.} \\
\texttt{Teine mängija saatis selle sõnumi:} \\
\\ 
\\ 
\texttt{"Kontrollisin oma laual (7,5) {-} see on tühi ◌. Sobib I{-}le. Liigutan I sinna ja sina liiguta ka oma I (7,5){-}sse. Alustan liigutamist."} \\
\\ 
\\ 
\texttt{Mis on sinu järgmine käsk?} \\
            }
        }
    }
    & \\ \\

    \theutterance \stepcounter{utterance}  
    & & & \multicolumn{4}{p{0.6\linewidth}}{
        \cellcolor[rgb]{0.9,0.9,0.9}{
            \makecell[{{p{\linewidth}}}]{
                \texttt{\tiny{[GM$\langle$P2]}}
                \texttt{LIIGUTA: I, (7, 5)} \\
            }
        }
    }
     \\ \\

    \theutterance \stepcounter{utterance}  
    & & & \multicolumn{2}{p{0.3\linewidth}}{
        \cellcolor[rgb]{0.9,0.9,0.9}{
            \makecell[{{p{\linewidth}}}]{
                \texttt{\tiny{[GM$|$GM]}}
                \texttt{LIIGUTA: I, (7, 5)} \\
            }
        }
    }
    & & \\ \\

    \theutterance \stepcounter{utterance}  
    & & & \multicolumn{2}{p{0.3\linewidth}}{
        \cellcolor[rgb]{0.9,0.9,0.9}{
            \makecell[{{p{\linewidth}}}]{
                \texttt{\tiny{[GM$|$GM]}}
                \texttt{I moved to (7, 5)} \\
            }
        }
    }
    & & \\ \\

    \theutterance \stepcounter{utterance}  
    & & \multicolumn{4}{p{0.6\linewidth}}{
        \cellcolor[rgb]{0.9,0.9,0.9}{
            \makecell[{{p{\linewidth}}}]{
                \texttt{\tiny{[P1$\langle$GM]}}
                \texttt{Sinu sõnum on teisele mängijale edastatud.} \\
\\ 
\\ 
\texttt{Sa mängid praegu vooru 6 maksimaalsest 20 voorust.} \\
\texttt{Teil on kollektiivselt kogunenud 3 karistust 12 karistusest.} \\
\texttt{Teine mängija liigutas objekti oma ruudustikul.} \\
\\ 
\\ 
\texttt{Mis on sinu järgmine käsk?} \\
            }
        }
    }
    & \\ \\

    \theutterance \stepcounter{utterance}  
    & \multicolumn{4}{p{0.6\linewidth}}{
        \cellcolor[rgb]{0.9,0.9,0.9}{
            \makecell[{{p{\linewidth}}}]{
                \texttt{\tiny{[P1$\rangle$GM]}}
                \texttt{LIIGUTA: I, (7, 5)} \\
            }
        }
    }
    & & \\ \\

    \theutterance \stepcounter{utterance}  
    & & & \multicolumn{2}{p{0.3\linewidth}}{
        \cellcolor[rgb]{0.9,0.9,0.9}{
            \makecell[{{p{\linewidth}}}]{
                \texttt{\tiny{[GM$|$GM]}}
                \texttt{LIIGUTA: I, (7, 5)} \\
            }
        }
    }
    & & \\ \\

    \theutterance \stepcounter{utterance}  
    & & & \multicolumn{2}{p{0.3\linewidth}}{
        \cellcolor[rgb]{0.9,0.9,0.9}{
            \makecell[{{p{\linewidth}}}]{
                \texttt{\tiny{[GM$|$GM]}}
                \texttt{I moved to (7, 5)} \\
            }
        }
    }
    & & \\ \\

    \theutterance \stepcounter{utterance}  
    & & \multicolumn{4}{p{0.6\linewidth}}{
        \cellcolor[rgb]{0.9,0.9,0.9}{
            \makecell[{{p{\linewidth}}}]{
                \texttt{\tiny{[GM$\rangle$P2]}}
                \texttt{Edukalt teisaldatud 'I' asukohta (7,5). Teie uuendatud ruudustik näeb välja selline:} \\
\texttt{```} \\
\texttt{ 1234567} \\
\texttt{╔═╤══╤══╗} \\
\texttt{║◌│◌◌│◌H║ 1} \\
\texttt{║◌└┬─┴──╢ 2} \\
\texttt{║◌◌│◌◌◌◌║ 3} \\
\texttt{╟─┬┘◌◌┌─╢ 4} \\
\texttt{║◌│WT◌│I║ 5} \\
\texttt{╟─┼─┬─┤◌║ 6} \\
\texttt{║◌│C│◌│◌║ 7} \\
\texttt{╚═╧═╧═╧═╝} \\
\texttt{```} \\
\texttt{Sa mängid praegu vooru 7 maksimaalsest 20 voorust.} \\
\texttt{Teil on kollektiivselt kogunenud 3 karistust 12 karistusest.} \\
\texttt{Teine mängija liigutas objekti oma ruudustikul.} \\
\\ 
\\ 
\texttt{Mis on sinu järgmine käsk?} \\
            }
        }
    }
    & \\ \\

    \theutterance \stepcounter{utterance}  
    & & & \multicolumn{4}{p{0.6\linewidth}}{
        \cellcolor[rgb]{0.9,0.9,0.9}{
            \makecell[{{p{\linewidth}}}]{
                \texttt{\tiny{[GM$\langle$P2]}}
                \texttt{Nüüd on I õiges kohas (7,5). Järgmisena pean liigutama H ja C. Minu H on (7,1) ja see peab minema (5,5). Minu C on (3,7) ja see peab minema (4,7). Kõigepealt liigutan H, sest (5,5) on tühi.} \\
\\ 
\texttt{LIIGUTA: H, (5, 5)} \\
            }
        }
    }
     \\ \\

    \theutterance \stepcounter{utterance}  
    & & & \multicolumn{2}{p{0.3\linewidth}}{
        \cellcolor[rgb]{0.9,0.9,0.9}{
            \makecell[{{p{\linewidth}}}]{
                \texttt{\tiny{[GM$|$GM]}}
                \texttt{Nüüd on I õiges kohas (7,5). Järgmisena pean liigutama H ja C. Minu H on (7,1) ja see peab minema (5,5). Minu C on (3,7) ja see peab minema (4,7). Kõigepealt liigutan H, sest (5,5) on tühi.} \\
\\ 
\texttt{LIIGUTA: H, (5, 5)} \\
            }
        }
    }
    & & \\ \\

    \theutterance \stepcounter{utterance}  
    & & & \multicolumn{2}{p{0.3\linewidth}}{
        \cellcolor[rgb]{0.9,0.9,0.9}{
            \makecell[{{p{\linewidth}}}]{
                \texttt{\tiny{[GM$|$GM]}}
                \texttt{Invalid format: head is not empty: 'Nüüd on I õiges kohas (7,5). Järgmisena pean liigutama H ja C. Minu H on (7,1) ja see peab minema (5,5). Minu C on (3,7) ja see peab minema (4,7). Kõigepealt liigutan H, sest (5,5) on tühi.} \\
\\ 
\texttt{'} \\
            }
        }
    }
    & & \\ \\

    \theutterance \stepcounter{utterance}  
    & & \multicolumn{4}{p{0.6\linewidth}}{
        \cellcolor[rgb]{0.9,0.9,0.9}{
            \makecell[{{p{\linewidth}}}]{
                \texttt{\tiny{[GM$\rangle$P2]}}
                \texttt{Karistus: Teie sõnum ei tohi sisaldada midagi enne käsku!} \\
\\ 
\texttt{Veendu, et sinu vastus sisaldab ainult kas `ÜTLE: <SÕNUM>` või `LIIGUTA: <OBJEKT>, (<X>, <Y>)` ja mitte midagi muud!} \\
\\ 
\texttt{Teil on kollektiivselt kogunenud 4 karistust 12 karistusest.} \\
\texttt{Palun proovi uuesti!} \\
            }
        }
    }
    & \\ \\

    \theutterance \stepcounter{utterance}  
    & & & \multicolumn{4}{p{0.6\linewidth}}{
        \cellcolor[rgb]{0.9,0.9,0.9}{
            \makecell[{{p{\linewidth}}}]{
                \texttt{\tiny{[GM$\langle$P2]}}
                \texttt{LIIGUTA: H, (5, 5)} \\
            }
        }
    }
     \\ \\

    \theutterance \stepcounter{utterance}  
    & & & \multicolumn{2}{p{0.3\linewidth}}{
        \cellcolor[rgb]{0.9,0.9,0.9}{
            \makecell[{{p{\linewidth}}}]{
                \texttt{\tiny{[GM$|$GM]}}
                \texttt{LIIGUTA: H, (5, 5)} \\
            }
        }
    }
    & & \\ \\

    \theutterance \stepcounter{utterance}  
    & & & \multicolumn{2}{p{0.3\linewidth}}{
        \cellcolor[rgb]{0.9,0.9,0.9}{
            \makecell[{{p{\linewidth}}}]{
                \texttt{\tiny{[GM$|$GM]}}
                \texttt{H moved to (5, 5)} \\
            }
        }
    }
    & & \\ \\

    \theutterance \stepcounter{utterance}  
    & & \multicolumn{4}{p{0.6\linewidth}}{
        \cellcolor[rgb]{0.9,0.9,0.9}{
            \makecell[{{p{\linewidth}}}]{
                \texttt{\tiny{[P1$\langle$GM]}}
                \texttt{Edukalt teisaldatud 'I' asukohta (7,5). Teie uuendatud ruudustik näeb välja selline:} \\
\texttt{```} \\
\texttt{ 1234567} \\
\texttt{╔═╤══╤══╗} \\
\texttt{║◌│◌◌│◌◌║ 1} \\
\texttt{║◌└┬─┴──╢ 2} \\
\texttt{║◌◌│◌◌◌◌║ 3} \\
\texttt{╟─┬┘◌◌┌─╢ 4} \\
\texttt{║◌│HW◌│I║ 5} \\
\texttt{╟─┼─┬─┤◌║ 6} \\
\texttt{║◌│C│T│◌║ 7} \\
\texttt{╚═╧═╧═╧═╝} \\
\texttt{```} \\
\texttt{Sa mängid praegu vooru 7 maksimaalsest 20 voorust.} \\
\texttt{Teil on kollektiivselt kogunenud 4 karistust 12 karistusest.} \\
\texttt{Teine mängija liigutas objekti oma ruudustikul.} \\
\\ 
\\ 
\texttt{Mis on sinu järgmine käsk?} \\
            }
        }
    }
    & \\ \\

    \theutterance \stepcounter{utterance}  
    & \multicolumn{4}{p{0.6\linewidth}}{
        \cellcolor[rgb]{0.9,0.9,0.9}{
            \makecell[{{p{\linewidth}}}]{
                \texttt{\tiny{[P1$\rangle$GM]}}
                \texttt{LIIGUTA: H, (5, 5)} \\
            }
        }
    }
    & & \\ \\

    \theutterance \stepcounter{utterance}  
    & & & \multicolumn{2}{p{0.3\linewidth}}{
        \cellcolor[rgb]{0.9,0.9,0.9}{
            \makecell[{{p{\linewidth}}}]{
                \texttt{\tiny{[GM$|$GM]}}
                \texttt{LIIGUTA: H, (5, 5)} \\
            }
        }
    }
    & & \\ \\

    \theutterance \stepcounter{utterance}  
    & & & \multicolumn{2}{p{0.3\linewidth}}{
        \cellcolor[rgb]{0.9,0.9,0.9}{
            \makecell[{{p{\linewidth}}}]{
                \texttt{\tiny{[GM$|$GM]}}
                \texttt{H moved to (5, 5)} \\
            }
        }
    }
    & & \\ \\

    \theutterance \stepcounter{utterance}  
    & & \multicolumn{4}{p{0.6\linewidth}}{
        \cellcolor[rgb]{0.9,0.9,0.9}{
            \makecell[{{p{\linewidth}}}]{
                \texttt{\tiny{[GM$\rangle$P2]}}
                \texttt{Edukalt teisaldatud 'H' asukohta (5,5). Teie uuendatud ruudustik näeb välja selline:} \\
\texttt{```} \\
\texttt{ 1234567} \\
\texttt{╔═╤══╤══╗} \\
\texttt{║◌│◌◌│◌◌║ 1} \\
\texttt{║◌└┬─┴──╢ 2} \\
\texttt{║◌◌│◌◌◌◌║ 3} \\
\texttt{╟─┬┘◌◌┌─╢ 4} \\
\texttt{║◌│WTH│I║ 5} \\
\texttt{╟─┼─┬─┤◌║ 6} \\
\texttt{║◌│C│◌│◌║ 7} \\
\texttt{╚═╧═╧═╧═╝} \\
\texttt{```} \\
\texttt{Sa mängid praegu vooru 8 maksimaalsest 20 voorust.} \\
\texttt{Teil on kollektiivselt kogunenud 4 karistust 12 karistusest.} \\
\texttt{Teine mängija liigutas objekti oma ruudustikul.} \\
\\ 
\\ 
\texttt{Mis on sinu järgmine käsk?} \\
            }
        }
    }
    & \\ \\

    \theutterance \stepcounter{utterance}  
    & & & \multicolumn{4}{p{0.6\linewidth}}{
        \cellcolor[rgb]{0.9,0.9,0.9}{
            \makecell[{{p{\linewidth}}}]{
                \texttt{\tiny{[GM$\langle$P2]}}
                \texttt{LIIGUTA: C, (4, 7)} \\
            }
        }
    }
     \\ \\

    \theutterance \stepcounter{utterance}  
    & & & \multicolumn{2}{p{0.3\linewidth}}{
        \cellcolor[rgb]{0.9,0.9,0.9}{
            \makecell[{{p{\linewidth}}}]{
                \texttt{\tiny{[GM$|$GM]}}
                \texttt{LIIGUTA: C, (4, 7)} \\
            }
        }
    }
    & & \\ \\

    \theutterance \stepcounter{utterance}  
    & & & \multicolumn{2}{p{0.3\linewidth}}{
        \cellcolor[rgb]{0.9,0.9,0.9}{
            \makecell[{{p{\linewidth}}}]{
                \texttt{\tiny{[GM$|$GM]}}
                \texttt{Karistus: (4,7) ei ole tühi, vaid sisaldab objekti '│'.} \\
\\ 
\texttt{Teil on kollektiivselt kogunenud 5 karistust 12 karistusest.} \\
\texttt{Palun proovi uuesti!} \\
            }
        }
    }
    & & \\ \\

    \theutterance \stepcounter{utterance}  
    & & \multicolumn{4}{p{0.6\linewidth}}{
        \cellcolor[rgb]{0.9,0.9,0.9}{
            \makecell[{{p{\linewidth}}}]{
                \texttt{\tiny{[GM$\rangle$P2]}}
                \texttt{Karistus: (4,7) ei ole tühi, vaid sisaldab objekti '│'.} \\
\\ 
\texttt{Teil on kollektiivselt kogunenud 5 karistust 12 karistusest.} \\
\texttt{Palun proovi uuesti!} \\
            }
        }
    }
    & \\ \\

    \theutterance \stepcounter{utterance}  
    & & & \multicolumn{4}{p{0.6\linewidth}}{
        \cellcolor[rgb]{0.9,0.9,0.9}{
            \makecell[{{p{\linewidth}}}]{
                \texttt{\tiny{[GM$\langle$P2]}}
                \texttt{LIIGUTA: C, (6, 7)} \\
            }
        }
    }
     \\ \\

    \theutterance \stepcounter{utterance}  
    & & & \multicolumn{2}{p{0.3\linewidth}}{
        \cellcolor[rgb]{0.9,0.9,0.9}{
            \makecell[{{p{\linewidth}}}]{
                \texttt{\tiny{[GM$|$GM]}}
                \texttt{LIIGUTA: C, (6, 7)} \\
            }
        }
    }
    & & \\ \\

    \theutterance \stepcounter{utterance}  
    & & & \multicolumn{2}{p{0.3\linewidth}}{
        \cellcolor[rgb]{0.9,0.9,0.9}{
            \makecell[{{p{\linewidth}}}]{
                \texttt{\tiny{[GM$|$GM]}}
                \texttt{Karistus: (6,7) ei ole tühi, vaid sisaldab objekti '│'.} \\
\\ 
\texttt{Teil on kollektiivselt kogunenud 6 karistust 12 karistusest.} \\
\texttt{Palun proovi uuesti!} \\
            }
        }
    }
    & & \\ \\

    \theutterance \stepcounter{utterance}  
    & & \multicolumn{4}{p{0.6\linewidth}}{
        \cellcolor[rgb]{0.9,0.9,0.9}{
            \makecell[{{p{\linewidth}}}]{
                \texttt{\tiny{[GM$\rangle$P2]}}
                \texttt{Karistus: (6,7) ei ole tühi, vaid sisaldab objekti '│'.} \\
\\ 
\texttt{Teil on kollektiivselt kogunenud 6 karistust 12 karistusest.} \\
\texttt{Palun proovi uuesti!} \\
            }
        }
    }
    & \\ \\

    \theutterance \stepcounter{utterance}  
    & & & \multicolumn{4}{p{0.6\linewidth}}{
        \cellcolor[rgb]{0.9,0.9,0.9}{
            \makecell[{{p{\linewidth}}}]{
                \texttt{\tiny{[GM$\langle$P2]}}
                \texttt{LIIGUTA: C, (7, 7)} \\
            }
        }
    }
     \\ \\

    \theutterance \stepcounter{utterance}  
    & & & \multicolumn{2}{p{0.3\linewidth}}{
        \cellcolor[rgb]{0.9,0.9,0.9}{
            \makecell[{{p{\linewidth}}}]{
                \texttt{\tiny{[GM$|$GM]}}
                \texttt{LIIGUTA: C, (7, 7)} \\
            }
        }
    }
    & & \\ \\

    \theutterance \stepcounter{utterance}  
    & & & \multicolumn{2}{p{0.3\linewidth}}{
        \cellcolor[rgb]{0.9,0.9,0.9}{
            \makecell[{{p{\linewidth}}}]{
                \texttt{\tiny{[GM$|$GM]}}
                \texttt{C moved to (7, 7)} \\
            }
        }
    }
    & & \\ \\

    \theutterance \stepcounter{utterance}  
    & & \multicolumn{4}{p{0.6\linewidth}}{
        \cellcolor[rgb]{0.9,0.9,0.9}{
            \makecell[{{p{\linewidth}}}]{
                \texttt{\tiny{[P1$\langle$GM]}}
                \texttt{Edukalt teisaldatud 'H' asukohta (5,5). Teie uuendatud ruudustik näeb välja selline:} \\
\texttt{```} \\
\texttt{ 1234567} \\
\texttt{╔═╤══╤══╗} \\
\texttt{║◌│◌◌│◌◌║ 1} \\
\texttt{║◌└┬─┴──╢ 2} \\
\texttt{║◌◌│◌◌◌◌║ 3} \\
\texttt{╟─┬┘◌◌┌─╢ 4} \\
\texttt{║◌│◌WH│I║ 5} \\
\texttt{╟─┼─┬─┤◌║ 6} \\
\texttt{║◌│C│T│◌║ 7} \\
\texttt{╚═╧═╧═╧═╝} \\
\texttt{```} \\
\texttt{Sa mängid praegu vooru 8 maksimaalsest 20 voorust.} \\
\texttt{Teil on kollektiivselt kogunenud 6 karistust 12 karistusest.} \\
\texttt{Teine mängija liigutas objekti oma ruudustikul.} \\
\\ 
\\ 
\texttt{Mis on sinu järgmine käsk?} \\
            }
        }
    }
    & \\ \\

    \theutterance \stepcounter{utterance}  
    & \multicolumn{4}{p{0.6\linewidth}}{
        \cellcolor[rgb]{0.9,0.9,0.9}{
            \makecell[{{p{\linewidth}}}]{
                \texttt{\tiny{[P1$\rangle$GM]}}
                \texttt{LIIGUTA: C, (4, 7)} \\
            }
        }
    }
    & & \\ \\

    \theutterance \stepcounter{utterance}  
    & & & \multicolumn{2}{p{0.3\linewidth}}{
        \cellcolor[rgb]{0.9,0.9,0.9}{
            \makecell[{{p{\linewidth}}}]{
                \texttt{\tiny{[GM$|$GM]}}
                \texttt{LIIGUTA: C, (4, 7)} \\
            }
        }
    }
    & & \\ \\

    \theutterance \stepcounter{utterance}  
    & & & \multicolumn{2}{p{0.3\linewidth}}{
        \cellcolor[rgb]{0.9,0.9,0.9}{
            \makecell[{{p{\linewidth}}}]{
                \texttt{\tiny{[GM$|$GM]}}
                \texttt{Karistus: (4,7) ei ole tühi, vaid sisaldab objekti '│'.} \\
\\ 
\texttt{Teil on kollektiivselt kogunenud 7 karistust 12 karistusest.} \\
\texttt{Palun proovi uuesti!} \\
            }
        }
    }
    & & \\ \\

    \theutterance \stepcounter{utterance}  
    & & \multicolumn{4}{p{0.6\linewidth}}{
        \cellcolor[rgb]{0.9,0.9,0.9}{
            \makecell[{{p{\linewidth}}}]{
                \texttt{\tiny{[P1$\langle$GM]}}
                \texttt{Karistus: (4,7) ei ole tühi, vaid sisaldab objekti '│'.} \\
\\ 
\texttt{Teil on kollektiivselt kogunenud 7 karistust 12 karistusest.} \\
\texttt{Palun proovi uuesti!} \\
            }
        }
    }
    & \\ \\

    \theutterance \stepcounter{utterance}  
    & \multicolumn{4}{p{0.6\linewidth}}{
        \cellcolor[rgb]{0.9,0.9,0.9}{
            \makecell[{{p{\linewidth}}}]{
                \texttt{\tiny{[P1$\rangle$GM]}}
                \texttt{LIIGUTA: T, (3, 5)} \\
            }
        }
    }
    & & \\ \\

    \theutterance \stepcounter{utterance}  
    & & & \multicolumn{2}{p{0.3\linewidth}}{
        \cellcolor[rgb]{0.9,0.9,0.9}{
            \makecell[{{p{\linewidth}}}]{
                \texttt{\tiny{[GM$|$GM]}}
                \texttt{LIIGUTA: T, (3, 5)} \\
            }
        }
    }
    & & \\ \\

    \theutterance \stepcounter{utterance}  
    & & & \multicolumn{2}{p{0.3\linewidth}}{
        \cellcolor[rgb]{0.9,0.9,0.9}{
            \makecell[{{p{\linewidth}}}]{
                \texttt{\tiny{[GM$|$GM]}}
                \texttt{T moved to (3, 5)} \\
            }
        }
    }
    & & \\ \\

    \theutterance \stepcounter{utterance}  
    & & \multicolumn{4}{p{0.6\linewidth}}{
        \cellcolor[rgb]{0.9,0.9,0.9}{
            \makecell[{{p{\linewidth}}}]{
                \texttt{\tiny{[GM$\rangle$P2]}}
                \texttt{Edukalt teisaldatud 'C' asukohta (7,7). Teie uuendatud ruudustik näeb välja selline:} \\
\texttt{```} \\
\texttt{ 1234567} \\
\texttt{╔═╤══╤══╗} \\
\texttt{║◌│◌◌│◌◌║ 1} \\
\texttt{║◌└┬─┴──╢ 2} \\
\texttt{║◌◌│◌◌◌◌║ 3} \\
\texttt{╟─┬┘◌◌┌─╢ 4} \\
\texttt{║◌│WTH│I║ 5} \\
\texttt{╟─┼─┬─┤◌║ 6} \\
\texttt{║◌│◌│◌│C║ 7} \\
\texttt{╚═╧═╧═╧═╝} \\
\texttt{```} \\
\texttt{Sa mängid praegu vooru 9 maksimaalsest 20 voorust.} \\
\texttt{Teil on kollektiivselt kogunenud 7 karistust 12 karistusest.} \\
\texttt{Teine mängija liigutas objekti oma ruudustikul.} \\
\\ 
\\ 
\texttt{Mis on sinu järgmine käsk?} \\
            }
        }
    }
    & \\ \\

    \theutterance \stepcounter{utterance}  
    & & & \multicolumn{4}{p{0.6\linewidth}}{
        \cellcolor[rgb]{0.9,0.9,0.9}{
            \makecell[{{p{\linewidth}}}]{
                \texttt{\tiny{[GM$\langle$P2]}}
                \texttt{ÜTLE: Liigutasin oma H (5,5) ja C (7,7). Aga C pidi minema (4,7), mis minu laual ei ole tühi {-} seal on joon. Kas (7,7) on sinu laual tühi? See on C uus asukoht minu laual. Minu objektid on nüüd: W (4,5), T (3,5), H (5,5), I (7,5), C (7,7). Kas need sobivad sinu lauale ühtimiseks?} \\
            }
        }
    }
     \\ \\

    \theutterance \stepcounter{utterance}  
    & & & \multicolumn{2}{p{0.3\linewidth}}{
        \cellcolor[rgb]{0.9,0.9,0.9}{
            \makecell[{{p{\linewidth}}}]{
                \texttt{\tiny{[GM$|$GM]}}
                \texttt{ÜTLE: Liigutasin oma H (5,5) ja C (7,7). Aga C pidi minema (4,7), mis minu laual ei ole tühi {-} seal on joon. Kas (7,7) on sinu laual tühi? See on C uus asukoht minu laual. Minu objektid on nüüd: W (4,5), T (3,5), H (5,5), I (7,5), C (7,7). Kas need sobivad sinu lauale ühtimiseks?} \\
            }
        }
    }
    & & \\ \\

    \theutterance \stepcounter{utterance}  
    & & \multicolumn{4}{p{0.6\linewidth}}{
        \cellcolor[rgb]{0.9,0.9,0.9}{
            \makecell[{{p{\linewidth}}}]{
                \texttt{\tiny{[P1$\langle$GM]}}
                \texttt{Edukalt teisaldatud 'T' asukohta (3,5). Teie uuendatud ruudustik näeb välja selline:} \\
\texttt{```} \\
\texttt{ 1234567} \\
\texttt{╔═╤══╤══╗} \\
\texttt{║◌│◌◌│◌◌║ 1} \\
\texttt{║◌└┬─┴──╢ 2} \\
\texttt{║◌◌│◌◌◌◌║ 3} \\
\texttt{╟─┬┘◌◌┌─╢ 4} \\
\texttt{║◌│TWH│I║ 5} \\
\texttt{╟─┼─┬─┤◌║ 6} \\
\texttt{║◌│C│◌│◌║ 7} \\
\texttt{╚═╧═╧═╧═╝} \\
\texttt{```} \\
\texttt{Sa mängid praegu vooru 9 maksimaalsest 20 voorust.} \\
\texttt{Teil on kollektiivselt kogunenud 7 karistust 12 karistusest.} \\
\texttt{Teine mängija saatis selle sõnumi:} \\
\\ 
\\ 
\texttt{"Liigutasin oma H (5,5) ja C (7,7). Aga C pidi minema (4,7), mis minu laual ei ole tühi {-} seal on joon. Kas (7,7) on sinu laual tühi? See on C uus asukoht minu laual. Minu objektid on nüüd: W (4,5), T (3,5), H (5,5), I (7,5), C (7,7). Kas need sobivad sinu lauale ühtimiseks?"} \\
\\ 
\\ 
\texttt{Mis on sinu järgmine käsk?} \\
            }
        }
    }
    & \\ \\

    \theutterance \stepcounter{utterance}  
    & \multicolumn{4}{p{0.6\linewidth}}{
        \cellcolor[rgb]{0.9,0.9,0.9}{
            \makecell[{{p{\linewidth}}}]{
                \texttt{\tiny{[P1$\rangle$GM]}}
                \texttt{ÜTLE: Kontrollisin (7,7) oma laual {-} see on tühi ◌. Aga (4,7) on minu laual '│', nii et C ei saa sinna minna. Sinu uus pakutud C asukoht (7,7) sobib mulle. Liigutan oma C (2,7) {-}> (7,7). Kinnita, et see on OK ja liiguta ka sina oma C (7,7){-}sse kui pole veel teinud.} \\
            }
        }
    }
    & & \\ \\

    \theutterance \stepcounter{utterance}  
    & & & \multicolumn{2}{p{0.3\linewidth}}{
        \cellcolor[rgb]{0.9,0.9,0.9}{
            \makecell[{{p{\linewidth}}}]{
                \texttt{\tiny{[GM$|$GM]}}
                \texttt{ÜTLE: Kontrollisin (7,7) oma laual {-} see on tühi ◌. Aga (4,7) on minu laual '│', nii et C ei saa sinna minna. Sinu uus pakutud C asukoht (7,7) sobib mulle. Liigutan oma C (2,7) {-}> (7,7). Kinnita, et see on OK ja liiguta ka sina oma C (7,7){-}sse kui pole veel teinud.} \\
            }
        }
    }
    & & \\ \\

    \theutterance \stepcounter{utterance}  
    & & \multicolumn{4}{p{0.6\linewidth}}{
        \cellcolor[rgb]{0.9,0.9,0.9}{
            \makecell[{{p{\linewidth}}}]{
                \texttt{\tiny{[GM$\rangle$P2]}}
                \texttt{Sinu sõnum on teisele mängijale edastatud.} \\
\\ 
\\ 
\texttt{Sa mängid praegu vooru 10 maksimaalsest 20 voorust.} \\
\texttt{Teil on kollektiivselt kogunenud 7 karistust 12 karistusest.} \\
\texttt{Teine mängija saatis selle sõnumi:} \\
\\ 
\\ 
\texttt{"Kontrollisin (7,7) oma laual {-} see on tühi ◌. Aga (4,7) on minu laual '│', nii et C ei saa sinna minna. Sinu uus pakutud C asukoht (7,7) sobib mulle. Liigutan oma C (2,7) {-}> (7,7). Kinnita, et see on OK ja liiguta ka sina oma C (7,7){-}sse kui pole veel teinud."} \\
\\ 
\\ 
\texttt{Mis on sinu järgmine käsk?} \\
            }
        }
    }
    & \\ \\

    \theutterance \stepcounter{utterance}  
    & & & \multicolumn{4}{p{0.6\linewidth}}{
        \cellcolor[rgb]{0.9,0.9,0.9}{
            \makecell[{{p{\linewidth}}}]{
                \texttt{\tiny{[GM$\langle$P2]}}
                \texttt{ÜTLE: Jah, minu C on juba (7,7) {-} liigutasin selle just sinna. Minu laual on nüüd kõik objektid: W (4,5), T (3,5), H (5,5), I (7,5), C (7,7). Kas sinu laual on kõik objektid nendes samades kohtades? Kui jah, siis oleme eesmärgi saavutanud ja võime mängu lõpetada. Kas oled valmis?} \\
            }
        }
    }
     \\ \\

    \theutterance \stepcounter{utterance}  
    & & & \multicolumn{2}{p{0.3\linewidth}}{
        \cellcolor[rgb]{0.9,0.9,0.9}{
            \makecell[{{p{\linewidth}}}]{
                \texttt{\tiny{[GM$|$GM]}}
                \texttt{ÜTLE: Jah, minu C on juba (7,7) {-} liigutasin selle just sinna. Minu laual on nüüd kõik objektid: W (4,5), T (3,5), H (5,5), I (7,5), C (7,7). Kas sinu laual on kõik objektid nendes samades kohtades? Kui jah, siis oleme eesmärgi saavutanud ja võime mängu lõpetada. Kas oled valmis?} \\
            }
        }
    }
    & & \\ \\

    \theutterance \stepcounter{utterance}  
    & & \multicolumn{4}{p{0.6\linewidth}}{
        \cellcolor[rgb]{0.9,0.9,0.9}{
            \makecell[{{p{\linewidth}}}]{
                \texttt{\tiny{[P1$\langle$GM]}}
                \texttt{Sinu sõnum on teisele mängijale edastatud.} \\
\\ 
\\ 
\texttt{Sa mängid praegu vooru 10 maksimaalsest 20 voorust.} \\
\texttt{Teil on kollektiivselt kogunenud 7 karistust 12 karistusest.} \\
\texttt{Teine mängija saatis selle sõnumi:} \\
\\ 
\\ 
\texttt{"Jah, minu C on juba (7,7) {-} liigutasin selle just sinna. Minu laual on nüüd kõik objektid: W (4,5), T (3,5), H (5,5), I (7,5), C (7,7). Kas sinu laual on kõik objektid nendes samades kohtades? Kui jah, siis oleme eesmärgi saavutanud ja võime mängu lõpetada. Kas oled valmis?"} \\
\\ 
\\ 
\texttt{Mis on sinu järgmine käsk?} \\
            }
        }
    }
    & \\ \\

    \theutterance \stepcounter{utterance}  
    & \multicolumn{4}{p{0.6\linewidth}}{
        \cellcolor[rgb]{0.9,0.9,0.9}{
            \makecell[{{p{\linewidth}}}]{
                \texttt{\tiny{[P1$\rangle$GM]}}
                \texttt{LIIGUTA: C, (7, 7)} \\
            }
        }
    }
    & & \\ \\

    \theutterance \stepcounter{utterance}  
    & & & \multicolumn{2}{p{0.3\linewidth}}{
        \cellcolor[rgb]{0.9,0.9,0.9}{
            \makecell[{{p{\linewidth}}}]{
                \texttt{\tiny{[GM$|$GM]}}
                \texttt{LIIGUTA: C, (7, 7)} \\
            }
        }
    }
    & & \\ \\

    \theutterance \stepcounter{utterance}  
    & & & \multicolumn{2}{p{0.3\linewidth}}{
        \cellcolor[rgb]{0.9,0.9,0.9}{
            \makecell[{{p{\linewidth}}}]{
                \texttt{\tiny{[GM$|$GM]}}
                \texttt{C moved to (7, 7)} \\
            }
        }
    }
    & & \\ \\

    \theutterance \stepcounter{utterance}  
    & & \multicolumn{4}{p{0.6\linewidth}}{
        \cellcolor[rgb]{0.9,0.9,0.9}{
            \makecell[{{p{\linewidth}}}]{
                \texttt{\tiny{[GM$\rangle$P2]}}
                \texttt{Sinu sõnum on teisele mängijale edastatud.} \\
\\ 
\\ 
\texttt{Sa mängid praegu vooru 11 maksimaalsest 20 voorust.} \\
\texttt{Teil on kollektiivselt kogunenud 7 karistust 12 karistusest.} \\
\texttt{Teine mängija liigutas objekti oma ruudustikul.} \\
\\ 
\\ 
\texttt{Mis on sinu järgmine käsk?} \\
            }
        }
    }
    & \\ \\

    \theutterance \stepcounter{utterance}  
    & & & \multicolumn{4}{p{0.6\linewidth}}{
        \cellcolor[rgb]{0.9,0.9,0.9}{
            \makecell[{{p{\linewidth}}}]{
                \texttt{\tiny{[GM$\langle$P2]}}
                \texttt{ÜTLE: valmis?} \\
            }
        }
    }
     \\ \\

    \theutterance \stepcounter{utterance}  
    & & & \multicolumn{2}{p{0.3\linewidth}}{
        \cellcolor[rgb]{0.9,0.9,0.9}{
            \makecell[{{p{\linewidth}}}]{
                \texttt{\tiny{[GM$|$GM]}}
                \texttt{ÜTLE: valmis?} \\
            }
        }
    }
    & & \\ \\

    \theutterance \stepcounter{utterance}  
    & & \multicolumn{4}{p{0.6\linewidth}}{
        \cellcolor[rgb]{0.9,0.9,0.9}{
            \makecell[{{p{\linewidth}}}]{
                \texttt{\tiny{[P1$\langle$GM]}}
                \texttt{Edukalt teisaldatud 'C' asukohta (7,7). Teie uuendatud ruudustik näeb välja selline:} \\
\texttt{```} \\
\texttt{ 1234567} \\
\texttt{╔═╤══╤══╗} \\
\texttt{║◌│◌◌│◌◌║ 1} \\
\texttt{║◌└┬─┴──╢ 2} \\
\texttt{║◌◌│◌◌◌◌║ 3} \\
\texttt{╟─┬┘◌◌┌─╢ 4} \\
\texttt{║◌│TWH│I║ 5} \\
\texttt{╟─┼─┬─┤◌║ 6} \\
\texttt{║◌│◌│◌│C║ 7} \\
\texttt{╚═╧═╧═╧═╝} \\
\texttt{```} \\
\texttt{Sa mängid praegu vooru 11 maksimaalsest 20 voorust.} \\
\texttt{Teil on kollektiivselt kogunenud 7 karistust 12 karistusest.} \\
\texttt{Teine mängija saatis selle sõnumi:} \\
\\ 
\\ 
\texttt{"valmis?"} \\
\\ 
\\ 
\texttt{Mis on sinu järgmine käsk?} \\
            }
        }
    }
    & \\ \\

    \theutterance \stepcounter{utterance}  
    & \multicolumn{4}{p{0.6\linewidth}}{
        \cellcolor[rgb]{0.9,0.9,0.9}{
            \makecell[{{p{\linewidth}}}]{
                \texttt{\tiny{[P1$\rangle$GM]}}
                \texttt{ÜTLE: valmis!} \\
            }
        }
    }
    & & \\ \\

    \theutterance \stepcounter{utterance}  
    & & & \multicolumn{2}{p{0.3\linewidth}}{
        \cellcolor[rgb]{0.9,0.9,0.9}{
            \makecell[{{p{\linewidth}}}]{
                \texttt{\tiny{[GM$|$GM]}}
                \texttt{ÜTLE: valmis!} \\
            }
        }
    }
    & & \\ \\

    \theutterance \stepcounter{utterance}  
    & & & \multicolumn{2}{p{0.3\linewidth}}{
        \cellcolor[rgb]{0.9,0.9,0.9}{
            \makecell[{{p{\linewidth}}}]{
                \texttt{\tiny{[GM$|$GM]}}
                \texttt{true} \\
            }
        }
    }
    & & \\ \\

    \theutterance \stepcounter{utterance}  
    & & & \multicolumn{2}{p{0.3\linewidth}}{
        \cellcolor[rgb]{0.9,0.9,0.9}{
            \makecell[{{p{\linewidth}}}]{
                \texttt{\tiny{[GM$|$GM]}}
                \texttt{Initial states:} \\
\texttt{```} \\
\texttt{Player 1:} \\
\texttt{ 1234567} \\
\texttt{╔═╤══╤══╗} \\
\texttt{║◌│◌◌│◌◌║ 1} \\
\texttt{║W└┬─┴──╢ 2} \\
\texttt{║◌◌│◌◌◌◌║ 3} \\
\texttt{╟─┬┘◌◌┌─╢ 4} \\
\texttt{║I│H◌◌│◌║ 5} \\
\texttt{╟─┼─┬─┤◌║ 6} \\
\texttt{║◌│C│T│◌║ 7} \\
\texttt{╚═╧═╧═╧═╝} \\
\\ 
\\ 
\texttt{Player 2:} \\
\texttt{ 1234567} \\
\texttt{╔═╤══╤══╗} \\
\texttt{║◌│◌◌│◌H║ 1} \\
\texttt{║◌└┬─┴──╢ 2} \\
\texttt{║◌◌│◌◌◌◌║ 3} \\
\texttt{╟─┬┘◌I┌─╢ 4} \\
\texttt{║◌│WT◌│◌║ 5} \\
\texttt{╟─┼─┬─┤◌║ 6} \\
\texttt{║◌│C│◌│◌║ 7} \\
\texttt{╚═╧═╧═╧═╝} \\
\\ 
\texttt{```} \\
            }
        }
    }
    & & \\ \\

    \theutterance \stepcounter{utterance}  
    & & & \multicolumn{2}{p{0.3\linewidth}}{
        \cellcolor[rgb]{0.9,0.9,0.9}{
            \makecell[{{p{\linewidth}}}]{
                \texttt{\tiny{[GM$|$GM]}}
                \texttt{End states:} \\
\texttt{```} \\
\texttt{Player 1:} \\
\texttt{ 1234567} \\
\texttt{╔═╤══╤══╗} \\
\texttt{║◌│◌◌│◌◌║ 1} \\
\texttt{║◌└┬─┴──╢ 2} \\
\texttt{║◌◌│◌◌◌◌║ 3} \\
\texttt{╟─┬┘◌◌┌─╢ 4} \\
\texttt{║◌│TWH│I║ 5} \\
\texttt{╟─┼─┬─┤◌║ 6} \\
\texttt{║◌│◌│◌│C║ 7} \\
\texttt{╚═╧═╧═╧═╝} \\
\\ 
\\ 
\texttt{Player 2:} \\
\texttt{ 1234567} \\
\texttt{╔═╤══╤══╗} \\
\texttt{║◌│◌◌│◌◌║ 1} \\
\texttt{║◌└┬─┴──╢ 2} \\
\texttt{║◌◌│◌◌◌◌║ 3} \\
\texttt{╟─┬┘◌◌┌─╢ 4} \\
\texttt{║◌│WTH│I║ 5} \\
\texttt{╟─┼─┬─┤◌║ 6} \\
\texttt{║◌│◌│◌│C║ 7} \\
\texttt{╚═╧═╧═╧═╝} \\
\\ 
\texttt{```} \\
            }
        }
    }
    & & \\ \\

    \theutterance \stepcounter{utterance}  
    & & & \multicolumn{2}{p{0.3\linewidth}}{
        \cellcolor[rgb]{0.9,0.9,0.9}{
            \makecell[{{p{\linewidth}}}]{
                \texttt{\tiny{[GM$|$GM]}}
                \texttt{* success: True} \\
\texttt{* lose: False} \\
\texttt{* aborted: False} \\
\texttt{{-}{-}{-}{-}{-}{-}{-}} \\
\texttt{* Shifts: 5.00} \\
\texttt{* Max Shifts: 8.00} \\
\texttt{* Min Shifts: 4.00} \\
\texttt{* End Distance Sum: 2.00} \\
\texttt{* Init Distance Sum: 15.62} \\
\texttt{* Expected Distance Sum: 20.95} \\
\texttt{* Penalties: 7.00} \\
\texttt{* Max Penalties: 12.00} \\
\texttt{* Rounds: 11.00} \\
\texttt{* Max Rounds: 20.00} \\
\texttt{* Object Count: 5.00} \\
            }
        }
    }
    & & \\ \\

\end{supertabular}
}

\end{document}
